\section{2024 Algebra \& Number Theory}

\subsection*{Exam}
\begin{exercise}
\label{algebra:2024:1}
Let $G$ be a finite group.

(1) Let $K$ be a field. Show that $G$ has a finite-dimensional faithful $K$-linear representation.

(2) Show that $G$ has a faithful one-dimensional complex representation if and only if $G$ is cyclic.

(3) Assume moreover that $G$ is commutative. Let $n \geq 1$ be an integer. Show that $G$ has a faithful $n$-dimensional complex representation if and only if $G$ can be generated by $n$ elements.

(4) Classify all finite groups having a faithful 2-dimensional real representation.
\end{exercise}

\begin{solution}
\textbf{Prerequisites:}
\begin{itemize}
    \item \textbf{Faithful Representation:} A representation $\rho: G \to \operatorname{GL}(V)$ is faithful if it is injective, meaning different group elements map to different matrices.
    \item \textbf{Regular Representation:} The action of a group on its own group algebra $K[G]$.
    \item \textbf{Cyclic Groups:} Groups generated by a single element. A key theorem states that any finite subgroup of the multiplicative group of a field is cyclic.
    \item \textbf{Character Theory:} For abelian groups, representations decompose into 1-dimensional characters.
\end{itemize}

\textbf{Steps:}

\textbf{(1) Faithful Representation over $K$:}
1. Define the vector space $V = K[G]$ with basis $\{e_g : g \in G\}$. The dimension is the order of the group $|G|$.
2. Define the action $\rho(g) e_h = e_{gh}$. This is called the left regular representation.
3. Check for faithfulness: If $\rho(g)$ acts as the identity, then $\rho(g) e_1 = e_1 \implies e_{g \cdot 1} = e_1 \implies g = 1$.
4. Since the kernel is trivial, the representation is faithful and finite-dimensional.

\textbf{(2) 1D Complex Representations:}
1. A 1D representation is a homomorphism $\chi: G \to \mathbb{C}^\times$.
2. If $\chi$ is faithful, $G$ is isomorphic to a finite subgroup of $\mathbb{C}^\times$. Since any finite subgroup of the multiplicative group of a field is cyclic, $G$ must be cyclic.
3. Conversely, if $G = \langle g \rangle$ is cyclic of order $n$, setting $\chi(g) = e^{2\pi i/n}$ provides a faithful character.

\textbf{(3) $n$-D Complex Representations for Commutative $G$:}
1. By the structure theorem, a finite abelian group $G$ is a product of cyclic groups.
2. If $G$ has a faithful $n$-D representation $\rho = \chi_1 \oplus \dots \oplus \chi_n$, then $G$ embeds into $(\mathbb{C}^\times)^n$. The image is a subgroup of a product of $n$ cyclic groups, hence can be generated by $n$ elements.
3. Conversely, if $G = \langle g_1, \dots, g_n \rangle$, we can take the direct sum of faithful characters of each cyclic factor to form an $n$-dimensional faithful representation.

\textbf{(4) 2D Real Representations:}
1. A finite subgroup of $\operatorname{GL}_2(\mathbb{R})$ is conjugate to a subgroup of the orthogonal group $O(2)$ (by averaging an inner product to make it $G$-invariant).
2. The finite subgroups of $O(2)$ are the cyclic groups $C_n$ (rotations) and the dihedral groups $D_{2n}$ (rotations and reflections).

\textbf{Intuition:}
Think of a representation as a way to "draw" the group as transformations of space. A faithful representation ensures that no information is lost—every group operation corresponds to a unique geometric movement. For 1D complex space, we are restricted to rotating on a circle, which only fits cyclic groups. In 2D real space, we can rotate and flip, which allows for the symmetries of regular polygons.
\end{solution}

\begin{exercise}
\label{algebra:2024:2}
Let $n \geq 1$ be an integer. Let $A$ be a DVR with $K$ its field of fractions and $\pi$ a uniformizer.
Show that $\mathbf{GL}_n(A) D_\mu \mathbf{GL}_n(A) \cap \mathbf{U}(K) D_\lambda$ is non-empty iff $\lambda_{\mathrm{dom}} \leq \mu_{\mathrm{dom}}$.
\end{exercise}

\begin{solution}
\textbf{Prerequisites:}
\begin{itemize}
    \item \textbf{DVR (Discrete Valuation Ring):} A ring like $\mathbb{Z}_p$ where divisibility by $\pi$ determines the "size" of elements.
    \item \textbf{Invariant Factors:} The powers of $\pi$ in the Smith Normal Form of a matrix, representing its intrinsic scale.
    \item \textbf{Dominance Order:} A way to compare sequences where $\lambda \leq \mu$ if the partial sums of $\lambda$ are smaller than those of $\mu$.
\end{itemize}

\textbf{Steps:}

1. \textbf{Identify the Sets:} $GL_n(A) D_\mu GL_n(A)$ is the set of matrices with invariant factors $\pi^{\mu_i}$. $\mathbf{U}(K) D_\lambda$ is the set of upper triangular matrices with diagonal entries $\pi^{\lambda_i}$.
2. \textbf{Relate Minors to Invariants:} A fundamental property of matrices over a DVR is that the valuation of any $k \times k$ minor must be at least the sum of the $k$ smallest valuations of the invariant factors.
3. \textbf{Calculate the Principal Minor:} For an upper triangular matrix $x$ with diagonal $\pi^{\lambda_i}$, the top-left $k \times k$ minor is exactly $\prod_{i=1}^k \pi^{\lambda_i}$, with valuation $\sum_{i=1}^k \lambda_i$.
4. \textbf{Apply the Bound:} This valuation $\sum_{i=1}^k \lambda_i$ must be $\geq \sum_{i=n-k+1}^n \mu_{\text{dom}, i}$ (sum of $k$ smallest invariants).
5. \textbf{Convert to Dominance:} Switching from smallest sums to largest sums (complementary sums), this condition is exactly equivalent to the dominance order $\lambda_{\text{dom}} \leq \mu_{\text{dom}}$.

\textbf{Intuition:}
The invariant factors $\mu$ represent the "true" sizes of the matrix's components. When you force the matrix into a triangular shape, the diagonal entries $\lambda$ become "averaged" or "smeared out" versions of these true sizes. The dominance order $\lambda \leq \mu$ mathematically captures this idea that $\lambda$ is "less spread out" than $\mu$.
\end{solution}

\begin{exercise}
\label{algebra:2024:3}
(1) $X^p - a$ is irreducible over imperfect $k$.
(2) Compute $A_{\mathrm{red}}$ for $A = k[X]/(X^{p^2} - aX^p)$.
\end{exercise}

\begin{solution}
\textbf{Prerequisites:}
\begin{itemize}
    \item \textbf{Imperfect Field:} A field of characteristic $p$ where $k \neq k^p$ (not every element is a $p$-th power).
    \item \textbf{Nilradical:} The set of nilpotent elements (elements $x$ such that $x^n=0$).
    \item \textbf{Reduced Ring ($A_{\text{red}}$):} The ring obtained by quotienting out the nilradical.
\end{itemize}

\textbf{Steps:}

\textbf{(1) Irreducibility:}
1. Suppose $X^p - a$ factors as $f(X)g(X)$. In the algebraic closure, $X^p - a = (X - a^{1/p})^p$.
2. Any factor $f(X)$ must be $(X - a^{1/p})^d$ for some $1 \leq d < p$.
3. The constant term of $f(X)$ is $\pm (a^{1/p})^d$, which must lie in the base field $k$.
4. Since $\gcd(d, p) = 1$, we can find $u, v$ such that $ud + vp = 1$.
5. Then $a^{1/p} = (a^{d/p})^u (a^{p/p})^v \in k$. This implies $a \in k^p$, contradicting that $k$ is imperfect. Thus $X^p - a$ is irreducible.

\textbf{(2) Computing $A_{\text{red}}$:}
1. Factor the polynomial: $X^{p^2} - aX^p = X^p(X^{p(p-1)} - a)$.
2. By the Chinese Remainder Theorem, $A \cong k[X]/(X^p) \times k[X]/(X^{p(p-1)} - a)$.
3. For the first factor, $X^p = 0$, so $X$ is nilpotent. Its reduced part is $k$.
4. For the second factor $R_2$, note that $(X^{p-1})^p = a$ in $R_2$. The element $X^{p-1} - a^{1/p}$ is nilpotent in the extension. The reduced part of $R_2$ is $k(a^{1/p})[X]/(X^{p-1} - a^{1/p})$.
5. This resulting polynomial is separable because its derivative is non-zero, so the ring is now reduced.
6. Thus $A_{\text{red}} \cong k \times k(a^{1/p})[X]/(X^{p-1} - a^{1/p})$.

\textbf{Intuition:}
Nilpotent elements are like "ghosts"—they appear in the equations but vanish when we focus on the core structure of the ring. Part (1) shows that if $a$ isn't a $p$-th power, $X^p-a$ is a solid, unbreakable block. Part (2) shows how we remove these "ghosts" (powers of $X$) to find the underlying field-like components of the ring.
\end{solution}

\begin{exercise}
\label{algebra:2024:4}
(1) $k((t))/k(t)$ is transcendental.
(2) Integral closure $B$ of $A = k[[t]] \cap K$ in $L$ is a DVR.
(3) $B$ is not finitely generated over $A$.
\end{exercise}

\begin{solution}
\textbf{Prerequisites:}
\begin{itemize}
    \item \textbf{$k((t))$ vs $k(t)$:} Power series (infinitely long) vs rational functions (finite fractions).
    \item \textbf{Integral Closure:} Elements in a larger field that satisfy monic polynomials over a subring.
    \item \textbf{Excellence:} A property of rings where closures behave nicely. Characteristic $p$ power series often lack this.
\end{itemize}

\textbf{Steps:}

\textbf{(1) Transcendence:}
1. $k(t)$ consists of rational functions, which are very simple. $k((t))$ contains series like $\sum t^{n!}$.
2. If $\sum t^{n!}$ were algebraic, it would satisfy a finite equation. However, the gaps between terms grow so fast that they cannot satisfy any polynomial relation.
3. (Or simply: the cardinality of $k((t))$ is much larger than $k(t)$).

\textbf{(2) $B$ is a DVR:}
1. $A$ and $B$ are intersections of the DVR $k[[t]]$ with certain subfields.
2. Any subring of a DVR that contains the uniformizer and is an intersection with a field is itself a DVR. The valuation is simply the power of $t$ in the series expansion.

\textbf{(3) Non-finite generation:}
1. If $B/A$ were a finite module, then because they have the same completion $k[[t]]$, the rank would have to be 1.
2. However, the degree of the field extension $[L:K]$ is $p$.
3. This discrepancy (rank 1 vs degree $p$) shows that $B$ cannot be finitely generated as an $A$-module.

\textbf{Intuition:}
Power series allow for "infinitely detailed" numbers. In characteristic $p$, the gap between "simple" numbers ($A$) and "detailed" numbers ($B$) can be so vast that you can't describe $B$ using just a finite number of pieces from $A$. This is why $B$ is not finitely generated.
\end{solution}

\begin{exercise}
\label{algebra:2024:5}
Let $p$ be a prime number. Consider $\mathbb{Z}_p$ (resp. $\mathbb{Q}_p$) the ring of $p$-adic integers (resp. field of $p$-adic numbers). Clearly $\mathbb{Z} \subset \mathbb{Z}_p$.

\begin{enumerate}
\item[(1)] Show that the set $\mathbb{Z}$ is dense in $\mathbb{Z}_p$, and deduce that a map $f: \mathbb{Z} \to \mathbb{Q}_p$ can be extended to a continuous function on $\mathbb{Z}_p$ if and only if $f$ is uniformly continuous.

\item[(2)] Let $a \in \mathbb{Q}_p \setminus \{0\}$. Under what condition on $a$, the map $f: \mathbb{Z} \to \mathbb{Q}_p, \quad n \mapsto a^n$ can be extended to a continuous function over $\mathbb{Z}_p$.

\item[(3)] Can we even extend $f$ to a continuous map $a^x: \mathbb{Q}_p \to \mathbb{Q}_p$ so that $a^{x+y} = a^x a^y$?
\end{enumerate}
\end{exercise}

\begin{solution}
\textbf{Prerequisites:}
\begin{itemize}
    \item \textbf{$p$-adic numbers:} A system where numbers are "close" if their difference is divisible by a high power of $p$.
    \item \textbf{Uniform Continuity:} The essential condition for extending a function from a dense set to its completion.
    \item \textbf{Divisible Group:} A group where you can always "divide" (e.g., in $\mathbb{Q}_p$, every $x$ has an $n$-th root).
\end{itemize}

\textbf{Steps:}

\textbf{(1) Density and Extension:}
1. $\mathbb{Z}_p$ is defined as the completion of $\mathbb{Z}$, so $\mathbb{Z}$ is dense by construction.
2. A standard analytical theorem states that a function on a dense set $D$ extends to the completion $\bar{D}$ if and only if it is uniformly continuous on $D$.

\textbf{(2) Extending $a^n$:}
1. For $a^n$ to be uniformly continuous, we need $a^n \to 1$ as $n$ gets $p$-adically small.
2. This requires $a$ to be "close to 1" in the $p$-adic sense. The condition is $|a-1|_p < 1$ (or $|a-1|_2 < 1/2$ for $p=2$).
3. This ensures that the sequence $a^{p^k}$ converges to 1.

\textbf{(3) Extension to $\mathbb{Q}_p$:}
1. \textbf{No.} If such an extension $a^x$ existed, it would be a continuous homomorphism.
2. The additive group $\mathbb{Q}_p$ is divisible, but the multiplicative group of values $1 + p\mathbb{Z}_p$ is not (e.g., it lacks $p$-th roots of many elements).
3. A homomorphism cannot map a divisible group to a non-divisible one.

\textbf{Intuition:}
The $p$-adic world rearranges numbers based on divisibility. To extend a function like $a^n$ to this new "line," the function must be very smooth (uniformly continuous). While we can stretch it to cover the $p$-adic integers, we can't stretch it to the whole field because the "multiplicative" geometry has holes that the "additive" geometry doesn't.
\end{solution}

\begin{exercise}
\label{algebra:2024:6}
(1) $\Phi_n$ mod $q$.
(2) $n = 2^{r+1}, p \equiv 5 \pmod 8$. Inner product, shortest vector.
\end{exercise}

\begin{solution}
\textbf{Prerequisites:}
\begin{itemize}
    \item \textbf{Cyclotomic Polynomial $\Phi_n$:} Its roots are primitive $n$-th roots of unity.
    \item \textbf{Frobenius Map ($x \mapsto x^q$):} Determines how roots are permuted in characteristic $q$.
    \item \textbf{Lattice:} A grid of points; the "shortest vector" is the non-zero point closest to zero.
\end{itemize}

\textbf{Steps:}

\textbf{(1) $\Phi_n$ mod $q$:}
1. The factors of $\Phi_n(X)$ mod $q$ correspond to the orbits of roots under the Frobenius $x \mapsto x^q$.
2. Each factor has degree $f$, where $f$ is the order of $q$ in the group $(\mathbb{Z}/n\mathbb{Z})^\times$.

\textbf{(2) Lattice Geometry:}
1. \textbf{Inner Product:} The lattice is $\mathbb{Z}[\zeta]$. The trace-based inner product $(\zeta^i, \zeta^j) = \delta_{ij} 2^r$ follows from the fact that the sum of roots of unity is zero unless the exponent is zero.
2. \textbf{Shortest Vector:} For $p \equiv 5 \pmod 8$, the prime $p$ splits into two ideals in the cyclotomic field.
3. For any $\alpha$ in such a prime ideal $\mathfrak{p}$, the norm $\text{Tr}(\alpha \bar{\alpha})$ is constrained by the ideal's properties.
4. Detailed calculation shows the shortest squared length is $2^r p$, so the shortest vector has length $\sqrt{2^r p}$.

\textbf{Intuition:}
In modular arithmetic, the Frobenius map acts like a "shuffler." The number of piles the roots end up in determines the number of factors. In the lattice part, we are placing points on a high-dimensional grid. The prime ideal condition acts like a "keep out" zone around the origin, forcing the grid points to be at least a certain distance ($\sqrt{2^r p}$) away.
\end{solution}
